\section{Τοπίο Αυτοματισμού σε Επίπεδο Αποθετηρίων}
\label{sec:appendix_workflow_landscape}

Στα \texttt{\ENG{scorm-engine}}, \texttt{\ENG{scorm-engine-php-sdk}},
\texttt{\ENG{example-lms-client}} και \texttt{\ENG{player}} εμφανίζονται σταθερά τα flows των \ENG{tests}, \ENG{lint},
\ENG{changelog}, \ENG{release} και \ENG{labeler}. Το γεγονός αυτό συνδέεται ισχυρά με την ύπαρξη του \texttt{\ENG{github-actions-template}},
το οποίο παράγει ακριβώς τέτοιου τύπου αρχεία για νέα αποθετήρια. Η ομοιομορφία αυτή έχει σημασία αλλά 
επίσης μας γλυτώνει από την ανάγκη να αναλύσουμε κάθε αποθετήριο ξεχωριστά, καθώς η βασική δομή των ροών είναι κοινή.

Η ύπαρξη \ENG{workflow automation} στο \texttt{\ENG{thesis-document-latex}} δεν πρέπει να συγχέεται με τον αυτοματισμό της πλατφόρμας.
Το αποθετήριο της διπλωματικής διαθέτει ξεχωριστά \ENG{LaTeX build}, \ENG{lint}, \ENG{spellcheck} και \ENG{release} workflows,
τα οποία τεκμηριώνουν μεν αναπαραγώγιμη παραγωγή του κειμένου, αλλά δεν συμμετέχουν στη λειτουργία του \ENG{SCORM ecosystem}.

\begin{figure}[H]
  \centering
\begin{chapterfigurebox}
  \begin{tikzpicture}[
      font=\small,
      repo/.style={draw, rounded corners, fill=gray!5, align=center, text width=2.9cm, minimum height=1.05cm},
      wf/.style={draw, rounded corners, fill=gray!12, align=center, text width=3.4cm, minimum height=1.05cm, font=\small\bfseries},
      arrow/.style={-Latex, thick}
    ]
    \node[repo] (engine) at (0,4.0) {\texttt{\ENG{scorm-engine}}};
    \node[repo] (sdk) at (0,2.4) {\texttt{\ENG{scorm-engine-php-sdk}}};
    \node[repo] (lms) at (0,0.8) {\texttt{\ENG{example-lms-client}}};
    \node[repo] (player) at (0,-0.8) {\texttt{\ENG{player}}};
    \node[repo] (infra) at (0,-2.4) {\texttt{\ENG{central-docker-infrastructure}}};

    \node[wf] (tests) at (5.6,3.2) {\ENG{tests + coverage}};
    \node[wf] (lint) at (5.6,1.2) {\ENG{lint / config validation}};
    \node[wf] (release) at (5.6,-0.8) {\ENG{changelog + release}};
    \node[wf] (integration) at (5.6,-2.8) {\ENG{compose integration}\\+\ENG{smoke checks}};

    \node[repo] (artifacts) at (11.0,2.2) {\ENG{coverage artifacts}\\\ENG{and Codecov}};
    \node[repo] (tags) at (11.0,0.2) {\ENG{CHANGELOGs, tags}\\\ENG{and releases}};
    \node[repo] (stack) at (11.0,-2.4) {\ENG{validated running stack}};

    \draw[arrow] (engine) -- (tests);
    \draw[arrow] (sdk) -- (tests);
    \draw[arrow] (lms) -- (tests);
    \draw[arrow] (player) -- (tests);
    \draw[arrow] (engine) -- (lint);
    \draw[arrow] (sdk) -- (lint);
    \draw[arrow] (lms) -- (lint);
    \draw[arrow] (player) -- (lint);
    \draw[arrow] (infra) -- (lint);
    \draw[arrow] (engine) -- (release);
    \draw[arrow] (sdk) -- (release);
    \draw[arrow] (lms) -- (release);
    \draw[arrow] (player) -- (release);
    \draw[arrow] (infra) -- (release);
    \draw[arrow] (infra) -- (integration);
    \draw[arrow] (tests) -- (artifacts);
    \draw[arrow] (release) -- (tags);
    \draw[arrow] (integration) -- (stack);
  \end{tikzpicture}
\end{chapterfigurebox}
  \caption{Συνοπτικός χάρτης των βασικών \ENG{workflow families} και των κύριων εξόδων τους σε επίπεδο οργανισμού.}
  \label{fig:appendix_cross_repo_workflow_map}
\end{figure}


\begin{table}[H]
\centering
\caption{Συγκριτική χαρτογράφηση των κύριων \ENG{GitHub Actions workflows} στα βασικά αποθετήρια}
\label{tab:appendix_repo_workflow_matrix}
\begingroup
\footnotesize
\setlength{\tabcolsep}{3pt}
\renewcommand{\arraystretch}{1.20}
\begin{adjustbox}{center,max width=\textwidth}
\begin{tabular}{|p{2.7cm}|p{3.4cm}|p{2.8cm}|p{3.6cm}|p{3.4cm}|}
\hline
Αποθετήριο & Κύριο \ENG{tests workflow} & \ENG{Lint/config workflow} & \ENG{Changelog/release workflow} & Σημαντική ιδιαιτερότητα \\
\hline
\texttt{\ENG{scorm-engine}} & \ENG{Maven tests} με \ENG{JaCoCo} & \ENG{Java lint/verify} & \ENG{semantic-release} και \ENG{conventional changelog} & \ENG{Coverage upload} προς \ENG{Codecov} υπό συνθήκη \\
\hline
\texttt{\ENG{scorm-engine-php-sdk}} & \ENG{PHPUnit} με \ENG{Clover} & \ENG{PHP syntax lint} και \ENG{composer validate} & \ENG{semantic-release} και \ENG{changelog workflow} & Στηρίζεται σε \ENG{PCOV} για \ENG{coverage generation} \\
\hline
\texttt{\ENG{example-lms-client}} & \ENG{Laravel tests} με \ENG{SQLite setup} & \ENG{PHP syntax lint} σε \ENG{Laravel paths} & \ENG{release} και αυτοματοποιημένο \ENG{changelog commit} & Χρησιμοποιεί \ENG{checkout} του \ENG{SDK} ως \ENG{path dependency} \\
\hline
\texttt{\ENG{player}} & \ENG{Node build/tests} με \ENG{c8} & \ENG{Node lint workflow} & \ENG{semantic-release} και \ENG{changelog workflow} & Συνδυάζει \ENG{build}, \ENG{unit tests} και \ENG{lcov export} \\
\hline
\texttt{\ENG{central-docker-infrastructure}} & Δεν υπάρχει \ENG{line-test suite} ανά υπηρεσία & \ENG{compose config}, \ENG{env} και \ENG{nginx validation} & Προσαρμοσμένο \ENG{tag/release workflow} & Διαθέτει ξεχωριστό \ENG{integration workflow} για την πλήρη σύνθεση \\
\hline
\end{tabular}
\end{adjustbox}
\endgroup

\end{table}
